\section{Introdução}

Os sistemas reator-separador possuem ampla relevância em diversos contextos industriais, sendo empregados em aplicações que vão da produção de polímeros, como o polietileno de baixa densidade \cite{hafele_2005}, a processos com reações exotérmicas e endotérmicas acopladas \cite{altimari_2008}.
Embora as etapas de reação e separação já interajam naturalmente nessas aplicações, ao retornar o reagente não convertido para o reator, em vez de descartá-lo, a recirculação viabiliza o aumento da conversão global em reações que sofrem limitação por equilíbrio.
Consequentemente, essa integração resulta na melhoria da seletividade do processo e na redução da necessidade de estágios adicionais de purificação \cite{sainio_2020}.

Entretanto, esse retorno de material dificulta a análise do sistema e intensifica o comportamento não linear inerente à cinética das reações, resultando em uma maior interação entre as variáveis do processo \cite{zeyer_2004, sainio_2020}.
A literatura reporta que esse nível de acoplamento pode induzir fenômenos atípicos, como a multiplicidade de estados estacionários, oscilações e até \textit{runaway} \cite{sainio_2020, bildea_2000, altimari_2008}, cujo desfecho depende diretamente da estratégia de operação adotada.
\textcite{pushpavanam_2001} demonstraram que fixar a vazão de alimentação fresca em uma rede reator-separador levou a estados instáveis e oscilações, enquanto a fixação da vazão de reciclo resultou em soluções estáveis para as mesmas condições.
Essa forte sensibilidade do sistema às escolhas operacionais e a possíveis distúrbios é um fator que contribui para agravar os desafios de controle e otimização na indústria.
Assim, o reciclo representa simultaneamente o mecanismo que melhora o desempenho de conversão e um fator que pode comprometer a estabilidade e a controlabilidade quando não é adequadamente analisado \cite{sainio_2020, zeyer_2004}.

Diante desse cenário, este trabalho tem como objetivo caracterizar o comportamento dinâmico de um sistema reator-separador acoplado por reciclo.
Especificamente, visa empregar ferramentas clássicas de análise de sistemas lineares para examinar o comportamento do processo nos domínios do tempo contínuo, da frequência e do tempo discreto.

% É essa dualidade entre relevância industrial e complexidade que motiva o presente estudo, cujo objetivo é caracterizar o comportamento transiente de um sistema reator-separador acoplado por reciclo. 

%Para isso, este trabalho estrutura-se na simulação do sistema para localizar seus pontos de equilíbrio, ao passo que o princípio da superposição é empregado para quantificar o grau de não linearidade do processo. A partir dessa etapa, empregam-se ferramentas clássicas de análise de sistemas lineares para examinar o comportamento do processo nos domínios do tempo contínuo, da frequência e do tempo discreto. Com isso, evidencia-se que as variações transitórias associadas à realimentação devem ser consideradas com a mesma atenção dada aos seus benefícios de conversão.
